\documentclass[book.tex]{subfiles}

\begin{document}
\chapter{Electrical Circuits}

\label{chap:electric_circuits}
\noindent\rule{11cm}{0.4pt} \\
\textbf{Prerequisites:} \autoref{chap:electric_current},  \\
\textbf{Difficulty Level:} * \\
\noindent\rule{11cm}{0.4pt}
\vspace{5mm}

An electrical circuit consists of individual electronic components such as resistors, capacitors, inductors, transistors and diode connected by conducting wires. Electrical circuits can be used to compute various functions on demand. This functionality means that electric circuits are tremendously useful for real physical devices. 

In this chapter, we introduce you to the basics components of electrical circuits and provide you with some mathematical tools to model them. We'll study both direct current and alternating current circuits as we make our tour.

\section{Basics of Circuits}

An electrical circuit is a mathematical abstraction of a real system with sources of electrical power (say a battery or other source) and various "circuit elements" which transform the input electrical power in various ways. Figure \ref{fig:circuit_diagram_and_schematic} demonstrates a drawing of a real circuit and its associated diagram.

\begin{figure}
    \centering
    % \includegraphics[width=0.7\linewidth]{figures/Physics/electromagnetism/electrical_circuits/circuit_diagram_and_schematic.png}
    \includegraphics[width=0.7\linewidth]{figures/Physics/electromagnetism/electrical_circuits/Real_Electrical_Circuit_1.png}
    \includegraphics[width=0.7\linewidth]{figures/Physics/electromagnetism/electrical_circuits/Real_Electrical_Circuit_2.png}
    \caption{Example of a real electrical circuit and its diagrammatic representation. }
    \label{fig:circuit_diagram_and_schematic}
\end{figure}
% \url{https://en.wikipedia.org/wiki/Circuit_diagram\#/media/File:Circuit_diagram_\%E2\%80\%93_pictorial_and_schematic.png}

Working with a diagrammatic abstraction of a circuit brings powerful advantages. Electrical engineers have developed a series of mathematical abstractions that will allow us to transform a given circuit into a set of equations which can be solved. As we will see later in this chapter, this same set of equations can be easily transformed into differentiable programs for modelling circuits.

\subsection{Basic Circuit Concepts}

We first need to introduce the basic physical concepts underpinning circuits. To start, we will stick to direct-current (DC) circuits in which the voltage does not vary or oscillate with time. Each current must have a voltage source which provides a fixed voltage difference $V$. The voltage difference will induce electrical charges to move and create an electrical current. We will refer to charge as $Q$ and current as $I$.

Resistors are a basic circuit element. Roughly put, a resistor is a region of the circuit where current can't freely flow. You can think of a resistor as an insulating element stuck into a circuit. Current can cross a resistive region, but at the cost of some energy. The fundamental equation which models a resistor is given by Ohm's law.

\section{Ohm's Law}
Suppose we have a circuit with a given resistor with resistance $R$. Then the voltage difference across the resistor $V$ is proportional to the product of the current $I$ moving across the resistor and the resistance.
\begin{equation}
    V = IR
\end{equation}

Let's turn this into a simple differentiable program.

\begin{lstlisting}
Voltage, Current, Resistance : $\mathbb{R}$
def ohms(I: Current, R: Resistance) $\to$ Voltage:
  return I * R
\end{lstlisting}
The next basic circuit element is the capacitor, which stores electric energy in the electric field. The capacitor has two separated plates. As charge $Q$ flows, a voltage difference $V$ builds up across the plates. The defining equation of a capacitor is given by

\begin{align}
C = \frac{Q}{V}
\end{align}

where $C$ is a quantity known as the capacitance.

Inductors are similar to capacitors but store energy in a magnetic field rather than in the electric field. The defining equation of an inductor is given by

\begin{align}
\Phi &= LI
\end{align}
where $L$ is a quantity called the inductance.

\subsection{Alternating Current Circuits}

Many real world systems have alternating current sources. Here, the voltage in the system is given by some time varying source $V(t)$. The current then becomes a time varying function $I(t)$ as well.

Impedence is the opposition a circuit presents to current when voltage is applied. The impedence is a complex number typically represented as $Z$. The impedence will prove a crucial quantity to analyze the behavior of an alternating current. Let's analyze the impedence of a simple alternating current circuit. We write the current as
\begin{align}
I(t)  &= |I|e^{j(\omega t + \phi_I)}
\end{align}
We write the voltage similarly as
\begin{align}
V(t) &= |V|e^{j(\omega t + \phi_V)}
\end{align}
The impedence is defined as the ratio of voltage to current
\begin{align}
Z &= \frac{V}{I} \\
&= \frac{|V|}{|I|}e^{j(\phi_V-\phi_I)}
\end{align}
Note that $Z$ does not vary in time. We can use these definitions to compute the impedence for resistors, capacitors, and inductors to be
\begin{align}
    Z_R &= R\\ 
    Z_C &= \frac{1}{j\omega C} \\
    Z_L &= j\omega L
\end{align}

\section{Transfer Functions}

For a continuous time input signal $x(t)$ and output signal $y(t)$, the transfer function $H(s)$ is the ratio of the Laplace transforms of the output and the input
\begin{align}
    H(s) = \frac{\mathcal{L}\{y(t)\}(s)}{\mathcal{L}\{x(t)\}(s)}
\end{align}
where $\mathcal{L}$ here denotes the Laplace transform and not a loss function. We write $s = \sigma + j\omega$ as a complex variable as the Laplace transformed variable. The transfer function proves a powerful tool to consider complex circuits with chained transformations.

%\textcolor{red}{TODO: Add section to math part on Laplace transform}

%\textcolor{red}{TODO: Add a section using the transfer function}

\section{Analyzing Circuits}

Mathematically analyzing a circuit can become complex quickly. In this section, we review a few standard methods for circuit analysis.

\subsection{Kirchhoff's Laws}

Kirchoff's laws provide standard rules for transforming an electrical circuit into a system of equations.

Kirchhoff's Circuit Law states that for any junction in an electrical circuit, the current entering the junction is equal to the current leaving the junction. Kirchhoff's Voltage Law states that the sum of the voltages around a closed loop is zero.

These two laws can be used to analyze complex circuits and transform them into equations which can be solved analytically and programmatically.

\subsection{RC Circuit}

Let's consider a simple resistor-capacitor circuit as in Figure \ref{rc_simple}. 

\begin{equation}
C \frac{dV}{dt} + \frac{V}{R} = 0
\end{equation}
where $C$ is the capacitance of the circuit and $V$ is the voltage across it. Let's solve this equation
\begin{align}
    C \frac{dV}{dt} + \frac{V}{R} &= 0 \\
    -\frac{1}{V} \frac{dV}{dt} &= \frac{1}{RC} \\
    - \int \frac{1}{V} \frac{dV}{dt} dt &= \frac{T}{RC} \\
    - \ln(V) - \ln(V_0) &= \frac{T}{RC} \\
    V &= V_0 e^{-\frac{CT}{R}}
\end{align}

Note how the voltage held in the capacitor dissipates over time.



\begin{marginfigure}
% \includegraphics[width=\textwidth]{figures/Physics/electromagnetism/electrical_circuits/RC_simple.png}
\includegraphics[width=\textwidth]{figures/Physics/electromagnetism/electrical_circuits/Simple_RC.png}
\caption{Simple RC Circuit. }
\label{rc_simple}
\end{marginfigure}
% TODO replace with our own image. Sourced from \url{https://en.wikipedia.org/wiki/RC_circuit\#/media/File:Discharging_capacitor.svg}

\begin{marginfigure}
% \includegraphics[width=\textwidth]{figures/Physics/electromagnetism/electrical_circuits/RC_circuit.png}
\includegraphics[width=\textwidth]{figures/Physics/electromagnetism/electrical_circuits/RC Circuit.png}
\caption{RC Circuit. }
\label{rc}
\end{marginfigure}
%TODO replace with our own image. Sourced from \url{https://en.wikipedia.org/wiki/Capacitor\#/media/File:RC_switch.svg}


%\textcolor{red}{TODO: Add analysis of RC circuit as a voltage transformation and consider its transfer function.}

%\textcolor{red}{TODO: Add physika code that models an RC circuit as a voltage transformation and optimize it to become a high pass or a low pass filter.}

\end{document}